\chapter{Dilation and Curettage (D\&C)}

\section*{Introduction \& Clinical Indications}
Dilation and curettage opens the cervical canal and removes tissue from the uterine cavity using suction, instruments, or both. It may be used for selected abnormal uterine bleeding, retained pregnancy tissue, miscarriage management, or tissue diagnosis. Ultrasound-guided aspiration, medication, hysteroscopy, or expectant management may be preferable depending on the indication. Blind curettage should not be presented as the default for all intrauterine pathology.

Pregnancy location and gestational age must be considered when pregnancy tissue is involved. A patient with pain, bleeding, syncope, or pregnancy of unknown location may require evaluation for ectopic pregnancy before D\&C. The procedure does not treat an ectopic pregnancy outside the uterus.

\section*{Relevant Surgical Anatomy}
The cervix connects the vagina to the uterine cavity. The uterine wall contains the endometrium, myometrium, and serosa; excessive instrumentation can perforate the myometrium or injure adjacent bowel, bladder, or vessels. Uterine position, size, fibroids, cesarean scars, cervical stenosis, and congenital anomalies affect access and risk.

\section*{Preoperative Workup \& Preparation}
Assessment includes history, examination, pregnancy testing, ultrasound when indicated, blood count, blood group and Rh status when clinically relevant, infection assessment, medication review, and anesthesia risk. Cervical preparation may reduce trauma. Consent includes pain, bleeding, infection, perforation, retained tissue, intrauterine adhesions, cervical injury, repeat procedure, anesthesia, and unexpected pathology.

\section*{Surgical Technique \& Procedure}
After anesthesia and positioning, the cervix is visualized and stabilized. It may be softened with medication or dilated gradually. Suction or curettage removes the planned tissue, with ultrasound guidance when it improves safety. Tissue is inspected and sent for pathology or other testing when indicated. Hemostasis and uterine integrity are assessed before discharge.

\section*{Complications \& Risks}
Risks include hemorrhage, infection, uterine perforation, cervical laceration, retained tissue, intrauterine adhesions, anesthesia complications, and need for re-aspiration or laparoscopy. Persistent pregnancy tissue, ongoing pregnancy, molar pregnancy, or ectopic pregnancy requires appropriate follow-up.

\section*{Postoperative Care \& Recovery}
Cramping and light bleeding may occur. The patient receives warning instructions for heavy bleeding, fever, severe pain, fainting, foul discharge, or persistent pregnancy symptoms. Pathology and hCG follow-up are arranged when relevant. Contraception can usually begin immediately, and ovulation may return rapidly.

\section*{Outcomes \& Prognosis}
Outcome depends on indication and completeness of evacuation. Most uncomplicated procedures have a short recovery, but fertility and menstrual effects should be discussed when adhesions or repeated procedures are possible. D\&C is not a substitute for hysteroscopic evaluation when focal intrauterine disease needs direct visualization.

\section*{References}
\begin{enumerate}
  \item American College of Obstetricians and Gynecologists. Early Pregnancy Loss: Practice Bulletin.
  \item American College of Obstetricians and Gynecologists. The Use of Hysteroscopy for Diagnosis and Treatment of Intrauterine Pathology.
  \item Royal College of Obstetricians and Gynaecologists. Best Practice in Abortion Care.
\end{enumerate}

\clearpage
